% Part: set-theory
% Chapter: z
% Section: powerset

\documentclass[../../../include/open-logic-section]{subfiles}

\begin{document}

\olfileid{sth}{z}{power}
\olsection{ঘাত সেট}

আরেকটি স্বতঃসিদ্ধ নিয়ে এগোব:

\begin{axiom}[ঘাত সেট]
যেকোনো সেট~$A$-এর জন্য
$\Pow{A} = \Setabs{x}{x \subseteq A}$ সেটটি আছে।
\[
  \forall A \exists P \forall x(x \in P \liff (\forall z \in x)z \in A)
\]
\end{axiom}

এর সমর্থনে যুক্তিটি বেশ সরল। ধরি, $A$ পর্যায়~$S$-এ
গঠিত হয়েছে। তাহলে $A$-এর সব সদস্য $S$-এর আগে
উপলব্ধ ছিল (\stagesacc{} অনুযায়ী)। তাই পৃথকীকরণের
সমর্থনে ব্যবহৃত যুক্তির মতোই, $A$-এর প্রত্যেক উপসেট
পর্যায়~$S$-এর মধ্যেই গঠিত হয়। ফলে $S$-এর পরের
যেকোনো পর্যায়ে সেগুলি সব উপলব্ধ থাকে, একটি
সেট হিসেবে গঠিত হওয়ার জন্য। এমন পর্যায় যে আছে
তা-ও জানি, কারণ $S$ শেষ পর্যায় নয়
(\stagessucc{} অনুযায়ী)। অতএব $\Pow{A}$ আছে।

ঘাত সেটের একটি সুন্দর ফল হলো:

\begin{prop}\label{thm:Products}
যেকোনো সেট $A, B$-এর কার্তেসীয় গুণফল
$A \times B$ আছে।
\end{prop}

\begin{proof}
ঘাত সেটের স্বতঃসিদ্ধ ও
\olref[sth][z][pairs]{prop:pairsconsequences} অনুযায়ী
$\Pow{\Pow{A \cup B}}$ আছে। তাই পৃথকীকরণ
অনুযায়ী এই সেটটিও আছে:
\[
  C = \Setabs{z \in \Pow{\Pow{A \cup B}}}{(\exists x \in A)(\exists y \in B) z = \tuple{x, y}}.
\]
এখন যেকোনো $x \in A$ ও~$y \in B$-এর জন্য
\olref[sth][z][pairs]{prop:pairsconsequences} অনুযায়ী
$\tuple{x, y}$ আছে। আরও, $x, y \in A \cup B$
হওয়ায় $\{x\}, \{x, y\} \in \Pow{A \cup B}$ এবং
$\tuple{x,y} \in \Pow{\Pow{A \cup B}}$।
তাই $A \times B = C$।
\end{proof}

এই প্রমাণে ঘাত সেটের স্বতঃসিদ্ধ পৃথকীকরণের সঙ্গে
কাজ করে। এতে বিস্ময়ের কিছু নেই। পৃথকীকরণ না
থাকলে ঘাত সেটের স্বতঃসিদ্ধ খুব বেশি \emph{শক্তিশালী}
হতো না। কারণ পৃথকীকরণই বলে দেয় একটি সেটের
কোন কোন উপসেটের অস্তিত্ব আছে; ফলে প্রতিটি ঘাত
সেট কতটা ``বড়'', তাও নির্ধারিত হয়।

\begin{prob}
দেখান যে যেকোনো সেট $A, B$-এর জন্য:
(i)~সংজ্ঞাক্ষেত্র~$A$ ও বিস্তৃতি~$B$ এমন সব সম্পর্কের
সেট আছে; এবং (ii)~$A$ থেকে~$B$-এ সব অপেক্ষকের
সেট আছে।
\end{prob}

\begin{prob}
ধরি, $A$ একটি সেট এবং $\sim$ $A$-এর উপর একটি
তুল্যতা সম্পর্ক। প্রমাণ করুন যে, $A$-এর উপর
$\sim$-এর সব তুল্যতা-শ্রেণির সেট, অর্থাৎ
$\equivclass{A}{\sim}$, আছে।
\end{prob}

\end{document}
